\documentclass[conference]{IEEEtran}
\usepackage[utf8]{inputenc}
\usepackage[T1]{fontenc}
\usepackage[spanish,es-tabla]{babel}
\usepackage{amsmath,booktabs,graphicx,cite}
\usepackage{microtype}
\hyphenation{aglu-ti-nan-te}
\renewcommand{\IEEEkeywordsname}{Palabras clave}
\begin{document}
\title{Traducción neuronal español–quechua ayacuchano con NLLB-200 ajustado sobre texto no religioso}
\author{
\IEEEauthorblockN{Flavio César Amayo Gamboa}
\IEEEauthorblockA{Maestría en Inteligencia Artificial \\ FIIS – UNI\\Lima, Perú\\flavio.amayo.g@uni.pe}
\and
\IEEEauthorblockN{Gustavo Adolfo Huarcaya Delgado}
\IEEEauthorblockA{Maestría en Inteligencia Artificial \\ FIIS – UNI\\Lima, Perú\\gustavo.huarcaya.d@uni.pe}
\and
\IEEEauthorblockN{Jhordy Joel Maximiliano Cusi}
\IEEEauthorblockA{Maestría en Inteligencia Artificial \\ FIIS – UNI\\Lima, Perú\\jmaximilianoc@uni.pe}
\and
\IEEEauthorblockN{Sergio Mondragon Contreras}
\IEEEauthorblockA{Maestría en Inteligencia Artificial \\ FIIS – UNI\\Lima, Perú\\sergio.mondragon.c@uni.pe}
}
\maketitle
\begin{abstract}
\textbf{[PENDIENTE: resumen con los resultados finales]}
\end{abstract}
\begin{IEEEkeywords}
traducción automática neuronal, quechua, lenguas de bajos recursos, NLLB, ajuste fino, chrF++
\end{IEEEkeywords}
\section{Introducción}
El quechua es una familia de lenguas con varios millones de hablantes en los Andes y es oficial en el Perú \cite{r3}. La variedad de Ayacucho (código ISO quy) sirve de base al quechua sureño estándar, que es el que usan las tareas compartidas de AmericasNLP \cite{r3,r4}. Es una lengua aglutinante: una palabra suele llevar varios sufijos de caso, persona y tiempo, de modo que el vocabulario observado crece rápido y cada forma aparece pocas veces. Hay además pocos textos paralelos con el español, y el recurso más usado, publicaciones de una sola organización, procede de publicaciones religiosas \cite{r11}.

Este trabajo aborda la traducción de oraciones en ambos sentidos, español a quechua y quechua a español. La entrada es una oración en un idioma y la salida es su traducción al otro. Medimos el éxito con chrF++ y BLEU sobre un conjunto de prueba que el modelo no vio ni en el entrenamiento ni en la selección de hiperparámetros.

Usamos aprendizaje profundo porque existe un modelo multilingüe preentrenado, NLLB-200, que ya incluye el quechua ayacuchano (quy\_Latn) \cite{r1}. Es un Transformer codificador-decodificador \cite{r2}, y ajustarlo con unas cien mil oraciones es viable en una sola GPU. No entrenamos desde cero porque unas cien mil oraciones no bastan para ello.

Los conjuntos públicos disponibles en Hugging Face tienen problemas que cambian los números reportados. En el conjunto de Ayacucho, el 7.3 \% de las oraciones en español de validación y el 6.9 \% de las de prueba ya estaban en entrenamiento. Más de la mitad de las filas son de contenido religioso, en buena parte del tipo de las publicaciones de publicaciones de una sola organización \cite{r11}. Y hay duplicados y pares mal alineados. Un traductor entrenado con ese material aprende sobre todo el estilo de esas publicaciones, y un promedio global sobre su conjunto de prueba no dice cuánto traduce fuera de él.

Nuestra decisión fue excluir el texto religioso del entrenamiento, la validación y la prueba, y entrenar solo con lo que queda: 49 414 pares de entrenamiento. Es menos de la mitad de los datos, así que medimos también cuánto aporta el tamaño del conjunto.

Las contribuciones son:

\begin{itemize}
\item Un conjunto de 122 802 pares español–quechua ayacuchano, limpiado y sin fuga entre particiones, y su subconjunto no religioso de 55 553 pares (49 414, 3 477 y 2 662 en entrenamiento, validación y prueba).
\item Un modelo bidireccional, NLLB-200 de 600 millones de parámetros ajustado con tres semillas y comparado con el mismo modelo sin ajustar.
\item Una ablación sobre el tamaño del entrenamiento (25 \%, 50 \% y 100 \% de los datos).
\item Un análisis de errores por longitud y sufijos, y el servicio y la aplicación Android que usan el modelo.
\end{itemize}
La Sección II resume trabajos previos, la III describe datos y método, la IV presenta los resultados, la V los discute y la VI concluye.

\section{Trabajos relacionados}
Sistemas para la tarea compartida de AmericasNLP. Moreno \cite{r3} ganó la pista sin datos de desarrollo en 2021: reunió corpus nuevos de poemas, canciones, léxicos y manuales, y preentrenó con español–inglés antes de ajustar para español–quechua. Su chrF de 34.6 es la referencia que cita el equipo de 2023 \cite{r4}. Ahmed et al. \cite{r4} añadieron textos literarios y legales y compararon un modelo preentrenado español–inglés con uno español–finés; el finés se eligió por ser aglutinante, y ese modelo rindió mejor, mientras que los datos sintéticos de baja calidad empeoraron los resultados. Asillo et al. \cite{r5} estudiaron si frases sintéticas generadas con un modelo de lenguaje grande, junto a un modelo preentrenado, mejoran la traducción del quechua.

Transferencia multilingüe. Chen y Abdul-Mageed \cite{r6} ajustaron modelos preentrenados bilingües y multilingües para traducir del español a diez lenguas indígenas sudamericanas, y fijaron el estado del arte en cinco de los diez pares sin ampliar los datos de entrenamiento. Mager et al. \cite{r7} y Tonja et al. \cite{r8} resumen el estado del área y sus carencias de datos.

Quechua a español. Ortega et al. \cite{r10} tradujeron del quechua al español con un modelo neuronal para una lengua polisintética y segmentación BPE guiada por sufijos. Cueva et al. \cite{r9} usaron un Transformer con LMfusion y segmentación PRPE, y construyeron una interfaz web de traducción.

Este trabajo se diferencia en tres puntos. Parte de una auditoría de los datos abiertos y reparte las particiones sin fuga. Excluye el texto religioso del entrenamiento y de la evaluación, y mide cuánto aporta el tamaño del conjunto que queda. Y entrega un modelo en ambos sentidos con servicio y aplicación, no solo una tabla de métricas.

\section{Metodología}
\subsection{Datos}
Usamos el conjunto spanish-to-quechua del hackathon SomosNLP 2022 en Hugging Face, variedad Ayacucho, con columnas es y qu. Reúne blogs, Tatoeba, datos de AmericasNLP 2021 y publicaciones de una sola organización \cite{r11}. No declara licencia; la parte que proviene de publicaciones de una sola organización tiene derechos de autor, y lo que queda tras el filtro puede seguir incluyendo texto de esa fuente, como se discute en la Sección V. Descartamos el conjunto de Cusco, también en Hugging Face, porque el quechua es una familia de variedades con inteligibilidad limitada entre las lejanas, en \cite{r4} añadir datos de Cusco ayudó menos que añadir datos de Ayacucho, y NLLB solo tiene la etiqueta de Ayacucho.

La limpieza une las particiones originales, aplica los filtros de la Tabla I y reparte de nuevo por grupos: todas las filas que comparten una oración en español o en quechua van a la misma partición, elegida con un hash del grupo. Después de esto, ninguna oración en español ni en quechua aparece en dos particiones. Cada fila se marca como religiosa o no. Un primer filtro de cinco palabras (Dios, Jehová, Biblia, Jesús, Cristo) marcó el 34.0 \% de las filas, pero una revisión de muestras mostró que dejaba pasar texto del mismo tipo: vocabulario de las revistas (congregación, predicar, adoración), citas como (Is. 55: 9), nombres bíblicos y formas quechuas con sufijos como Diospa. El filtro final amplía la lista con ese vocabulario y con patrones de cita, y busca en ambos idiomas. Es conservador: ante la duda marca la fila como religiosa, y por eso descarta también algunas oraciones generales con palabras como alma o Padre, y puede dejar algo de texto religioso sin detectar. Marca el 55.0 \% del entrenamiento, el 49.4 \% de la validación y el 56.5 \% de la prueba. Esas filas se excluyen de todas las particiones.

\begin{table}[!t]
\caption{Resumen de la limpieza (variedad Ayacucho)}
\label{tab:limpieza}
\centering\small
\begin{tabular}{lc}
\toprule
Paso & Filas \\
\midrule
Filas brutas (dos particiones originales unidas) & 128 434 \\
Vacía o de menos de 3 caracteres & 62 \\
Español idéntico al quechua & 22 \\
Más de 400 caracteres & 77 \\
Razón de longitud mayor a 4 & 605 \\
Sin letras & 316 \\
Duplicado exacto & 4 550 \\
Filas limpias & 122 802 \\
\bottomrule
\end{tabular}
\end{table}
Antes de excluir el texto religioso las particiones tenían 109 807, 6 877 y 6 118 pares. Después quedan 49 414 de entrenamiento, 3 477 de validación y 2 662 de prueba.

\subsection{Modelo}
Partimos de NLLB-200-distilled-600M \cite{r1}, un Transformer \cite{r2} de codificador y decodificador con etiquetas de idioma en la entrada. Un solo modelo traduce en ambos sentidos: cada par genera dos ejemplos, uno de es a qu (etiquetas spa\_Latn y quy\_Latn) y otro de qu a es. Todos los parámetros se ajustan; no se congela ninguna capa ni se añaden adaptadores.

\subsection{Función de pérdida y entrenamiento}
La pérdida es la entropía cruzada con forzado del maestro. Para un lote B de pares (x, y), con x la oración fuente e y la traducción de T tokens:

\begin{equation}
\mathcal{L}(\theta) = -\frac{1}{|B|}\sum_{(x,y)\in B}\frac{1}{T}\sum_{t=1}^{T}\log p_\theta\left(y_t \mid y_{<t}, x\right).
\end{equation}
Optimizamos con AdamW \cite{r14}, tasa de aprendizaje $3\cdot 10^{-5}$ y decaimiento de pesos 0.01, calentamiento lineal del 5 \% de los pasos y decaimiento lineal después, recorte de gradiente a norma 1, lotes de 32 oraciones de una sola dirección y 3 épocas. La longitud máxima es de 128 tokens. Los pesos se mantienen en fp32 y el cálculo usa bf16 con autocast (fp16 con escalado de gradiente en GPU sin bf16, como la T4 de Colab). Cada época evalúa 500 oraciones de validación no religiosas con chrF++ medio de ambas direcciones, y se conservan los pesos de la mejor época. La prueba se evalúa una sola vez con esos pesos, con búsqueda en haz de ancho 4. Las semillas son 0, 1 y 2; cambian el orden de los lotes y las 500 oraciones de validación. La ablación usa una fracción aleatoria del entrenamiento (25 \% y 50 \%) con la semilla 0.

\subsection{Línea base y métricas}
La línea base es NLLB-200-distilled-600M sin ajustar (zero-shot) sobre el mismo conjunto de prueba sin texto religioso. Como referencia publicada, el sistema ganador de 2021 obtuvo 34.6 de chrF en la prueba de AmericasNLP \cite{r3,r4}; no es comparable de forma directa porque el conjunto de prueba es otro. La métrica principal es chrF++ \cite{r12}, más adecuada que BLEU para una lengua aglutinante, y reportamos también BLEU de sacreBLEU \cite{r13}. Los resultados se separan por dirección y se promedian entre semillas con su desviación estándar.

\section{Resultados}
\textbf{[PENDIENTE: los experimentos siguen en ejecución; esta sección se completa con las tablas y figuras finales]}

\section{Discusión}
\textbf{[PENDIENTE: se redacta con los resultados]}

\section{Conclusiones}
\textbf{[PENDIENTE: se redacta con los resultados]}

\section*{Contribuciones}
\begin{tabular}{@{}p{0.32\columnwidth}p{0.6\columnwidth}@{}}
\toprule
Integrante & Tareas en las que participó \\
\midrule
Flavio César Amayo Gamboa & (por completar por el equipo) \\
Gustavo Adolfo Huarcaya Delgado & (por completar por el equipo) \\
Jhordy Joel Maximiliano Cusi & (por completar por el equipo) \\
Sergio Mondragon Contreras & (por completar por el equipo) \\
\bottomrule
\end{tabular}

\section*{Disponibilidad de código y uso de herramientas de IA}
Código, configuraciones y resultados: repositorio del proyecto (etiqueta v1.0-final). Los pesos del mejor modelo se publican con el enlace indicado en el README. El modelo base NLLB-200 tiene licencia CC-BY-NC 4.0, por lo que el modelo ajustado se limita a uso no comercial. Se usó Claude Code (Anthropic) para escribir y depurar los scripts de limpieza, entrenamiento y análisis, el servicio, la aplicación Android y los borradores de redacción. El equipo ejecutó los experimentos y revisó los textos; todas las cifras provienen de esas ejecuciones.

\begin{thebibliography}{99}
\bibitem{r1} NLLB Team, "No language left behind: Scaling human-centered machine translation," arXiv:2207.04672, 2022.
\bibitem{r2} A. Vaswani, N. Shazeer, N. Parmar, J. Uszkoreit, L. Jones, A. N. Gomez, Ł. Kaiser, and I. Polosukhin, "Attention is all you need," in Advances in Neural Information Processing Systems (NeurIPS), vol. 30, 2017.
\bibitem{r3} O. Moreno, "The REPUcs' Spanish–Quechua submission to the AmericasNLP 2021 shared task on open machine translation," in Proc. 1st Workshop on NLP for Indigenous Languages of the Americas, 2021, pp. 241–247.
\bibitem{r4} N. Ahmed, N. Flechas Manrique, and A. Petrović, "Enhancing Spanish-Quechua machine translation with pre-trained models and diverse data sources: LCT-EHU at AmericasNLP shared task," in Proc. AmericasNLP Workshop, 2023, pp. 156–162.
\bibitem{r5} J. Asillo Congora, J. Santisteban, and R. Lazo Vasquez, "UCSP submission to the AmericasNLP 2025 shared task," in Proc. 5th Workshop on NLP for Indigenous Languages of the Americas, 2025, pp. 84–91.
\bibitem{r6} W.-R. Chen and M. Abdul-Mageed, "Improving neural machine translation of indigenous languages with multilingual transfer learning," arXiv:2205.06993, 2022.
\bibitem{r7} M. Mager, R. Bhatnagar, G. Neubig, N. T. Vu, and K. Kann, "Neural machine translation for the indigenous languages of the Americas: An introduction," arXiv:2306.06804, 2023.
\bibitem{r8} A. L. Tonja, F. Balouchzahi, S. Butt, O. Kolesnikova, H. Ceballos, A. Gelbukh, and T. Solorio, "NLP progress in indigenous Latin American languages," arXiv:2404.05365, 2024.
\bibitem{r9} B. Cueva Medina, G. F. Tuco Casquino, and J. Sulla-Torres, "Development of a neural machine translation model optimized with BERT for translation from Quechua to Spanish," in 22nd LACCEI Int. Multi-Conf. for Engineering, Education, and Technology, 2024.
\bibitem{r10} J. E. Ortega, R. Castro Mamani, and K. Cho, "Neural machine translation with a polysynthetic low resource language," Machine Translation, 2020, doi: 10.1007/s10590-020-09255-9.
\bibitem{r11} Ž. Agić and I. Vulić, "publicaciones de una sola organización: A wide-coverage parallel corpus for low-resource languages," in Proc. 57th Annual Meeting of the ACL, 2019.
\bibitem{r12} M. Popović, "chrF++: words helping character n-grams," in Proc. 2nd Conf. on Machine Translation (WMT), 2017, pp. 612–618.
\bibitem{r13} M. Post, "A call for clarity in reporting BLEU scores," in Proc. 3rd Conf. on Machine Translation (WMT), 2018, pp. 186–191.
\bibitem{r14} I. Loshchilov and F. Hutter, "Decoupled weight decay regularization," in Proc. Int. Conf. on Learning Representations (ICLR), 2019.
\end{thebibliography}
\end{document}
